\scanpage{31}
% archive sheet number: 30
% L1--7
Il semble néanmoins que dans le cas général, pour prouver qu’un fibré principal $E$ sur $X-S$ provient d’une \deleted{fibré principal} extension de $\unsure{\mathscr P|\mathcal P}_S$ par $G$, il suffit (utilisant la théorie de la descente) de le vérifier en \unsure{remplaçant} $S$ par les spectres des \uline{complétés} de ses anneaux locaux.
% L8--10
De plus, utilisant la géométrie formelle [lorsque $G$ est bien plongé dans le groupe \deleted{\unsure{H|\mathrm H}} linéaire] on est ramené à le vérifier \uline{formellement}.
% L11--14
Cela revient au cas d’un \unsure{anneau} de base \unsure{Artinien}. Travaillant de proche en proche, on est ramené à le vérifier \unsure{pour les courbes} \uline{\unsure{sur}} les corps résiduels de $S$.
% L15--19
Utilisant à nouveau la descente on serait amené à le vérifier à nouveau sur les clôtures algébriques des corps résiduels\unsure{...---} On peut sans doute s’en tirer plus \unsure{astucieusement}.
